\chapter*{Abstract} \label{sec:03-abstract}
\addcontentsline{toc}{chapter}{Abstract}

\begin{foldable}
  Recent advances in \acl{AI} have been driven by increasingly large and complex models, trained on increasingly large datasets.
  Yet, model and dataset size are practical bottlenecks across all stages of a model's life cycle, from data preparation and training through to deployment and storage.
  Among compression techniques for both models and datasets, \ac{KD} stands out because it treats compression as a knowledge transfer problem rather than a reduction problem.
  Specifically, model distillation trains a small student model to mimic the outputs of a larger teacher model.
  Dataset distillation, in turn, synthesizes a small, informative dataset that yields downstream performance comparable to its larger counterpart.
\end{foldable}

\begin{foldable}
  However, several practical constraints limit the applicability of \ac{KD} in real-world settings.
  Model distillation methods typically assume full access to the original training dataset and to the internal parameters of the teacher network.
  This assumption routinely fails in practice, where data is limited or unavailable because of proprietary or privacy restrictions, and where teacher models are exposed only as black-box interfaces.
  Dataset distillation, on the other hand, has matured for image domains but remains underexplored for text, where the discrete and compositional nature of language poses distinct challenges.
  The distilled text dataset must also withstand extreme compression ratios while remaining human-readable for auditability and interpretability.
\end{foldable}

\begin{foldable}
  This thesis addresses \ac{KD} under three restrictive settings: (i) few-shot access to real data with a black-box teacher, (ii) complete absence of real data with a black-box teacher, and (iii) text dataset distillation, where the goal is to synthesize a tiny surrogate dataset that preserves model training utility.
  We motivate our approaches through the lens of \textit{distributional fidelity}: whether the compact artefact, the student model or the surrogate corpus, faithfully represents the original.
  For the first setting, we propose \textit{\acl{DivBFKD}} (\ac{DivBFKD}), a framework that improves the diversity of the limited training data by generating high-quality synthetic samples, supporting black-box \ac{KD} from only a few samples per class.
  For the second setting, we propose \textit{\acl{DIPKD}} (\ac{DIPKD}), a data-free black-box \ac{KD} method that synthesizes training images by incorporating structured image priors, then optimizes them for diversity via a contrastive objective.
  For the third setting, we propose \textit{\acl{TAKE}} (\ac{TAKE}), which distills a large corpus into a compact dataset of high-influence, human-readable synthetic texts.
  Across all three settings, our methods reach state-of-the-art performance on \ac{KD} benchmarks, showing that effective distillation is achievable under severe practical constraints.
\end{foldable}
